\documentclass[12pt]{article}
\usepackage[utf8]{inputenc}
\usepackage[T1]{fontenc}
\usepackage[margin=1in]{geometry}
\usepackage{setspace}
\usepackage{xcolor}
\usepackage[colorlinks=true, urlcolor=blue!60!black, linkcolor=blue!60!black]{hyperref}
\setlength{\parindent}{0em}
\setlength{\parskip}{0.6em}
\renewcommand{\baselinestretch}{1.0}
\begin{document}
\begin{center}
{\LARGE \textit{Martin v. Hunter's Lessee}, 14 U.S. (1 Wheat.) 304 (1816)\par}
\vspace{0.5em}
{\large Case Brief --- Two-Page Edition\par}
\end{center}

\section{Memory Jogger}
The Supreme Court may review state-court judgments that decide federal questions; Section~25 of the Judiciary Act of 1789 is constitutional (\href{https://supreme.justia.com/cases/federal/us/14/304/}{Justia}).

\section{Detailed Case Facts}
Lord Fairfax devised Northern Neck lands to his British nephew Denny Fairfax; Virginia later granted a parcel to David Hunter. The 1783 peace treaty and 1794 Jay Treaty protected British subjects' lands. Chief Justice Marshall, who had a financial interest in the estate, recused.

\section{Procedural History}
The Supreme Court in 1813 upheld Fairfax's title under the Jay Treaty and sent a mandate. The Virginia Court of Appeals, led by Spencer Roane, refused to obey, holding Section~25 unconstitutional as applied to state courts. Martin's heir brought a second writ of error.

\section{Judicial Votes}
Reversed, 6--0. Story for the Court (Washington, Livingston, Todd, Duvall); Johnson concurring; Marshall recused.

\section{Holding}
Article~III appellate jurisdiction reaches state-court judgments on federal questions because jurisdiction attaches to the case, not the court. Section~25 is valid, a refusal to obey a mandate is reviewable, and the Court may examine intertwined state-law title.

\section{Analysis of Opinions}
\textbf{Story:} The people, not the States, ordained the Constitution. Article~III's appellate grant has no limit on the tribunal appealed from. The Supremacy Clause contemplates state judges deciding federal questions, local bias may sway them, and uniformity requires one revising court; the First Congress agreed. His claim that Congress must vest all first-tier jurisdiction was dicta. \textbf{Johnson:} the Court acts ``over persons and cases,'' not over state courts as subordinates.

\section{Examples: Future Applications}
\textbf{Same-Side:} A state criminal conviction rejecting a federal defense is reviewable (\textit{Cohens v. Virginia}). \textbf{Opposite-Side:} A judgment resting on an adequate and independent state ground is not. \textbf{Fence-Sitter:} Stripping all federal jurisdiction over a class of constitutional claims pits \textit{Martin}'s narrow holding against Story's mandatory-vesting dicta.

\section{Critique}
Roane called federal review a ``commandeering'' of state courts; modern originalists accept the holding but split over Story's dicta. Scholars treat \textit{Martin} as the linchpin of national rights, though it helped build judicial supremacy.

\section{Key Quotations}
\begin{quote}
``That motive is the importance, and even necessity, of uniformity of decisions throughout the whole United States upon all subjects within the purview of the Constitution. Judges of equal learning and integrity in different States might differently interpret a statute or a treaty of the United States, or even the Constitution itself; if there were no revising authority to control these jarring and discordant judgments and harmonize them into uniformity, the laws, the treaties, and the Constitution of the United States would be different in different States, and might perhaps never have precisely the same construction, obligation, or efficacy in any two States.''
\end{quote}
\noindent --- Justice Story, Opinion of the Court, 14 U.S. (1 Wheat.) at 347--48.

\end{document}
